% Vorzeichen des Levi-Civita-Symbols: zyklisch +1, antizyklisch -1
\input{style}
\begin{document}
\begin{tikzpicture}
  % links: zyklische Reihenfolge 1-2-3, positives Vorzeichen
  \begin{scope}
    \foreach \a/\n in {90/1,210/2,330/3} {
      \node[circle, fill=fg, inner sep=0pt, minimum size=4pt] at (\a:1.25) {};
      \node at (\a:1.75) {$\n$};
    }
    % Bögen an festen Winkeln statt shorten: so liegt die Spitze tangential auf dem Bogen
    \foreach \s in {90,210,330}
      \draw[vec, acc, -{Stealth[length=6pt,width=5pt,bend]}]
        (\s+17:1.25) arc[start angle=\s+17, delta angle=86, radius=1.25];
    \node[acc] at (0,0) {$+1$};
    \node[anchor=north, font=\footnotesize] at (0,-2.1)
      {zyklisch: $123,\ 231,\ 312$};
  \end{scope}

  % rechts: gleiche Lage, Umlauf in Gegenrichtung 1-3-2, negatives Vorzeichen
  \begin{scope}[xshift=4.8cm]
    \foreach \a/\n in {90/1,210/2,330/3} {
      \node[circle, fill=fg, inner sep=0pt, minimum size=4pt] at (\a:1.25) {};
      \node at (\a:1.75) {$\n$};
    }
    \foreach \s in {90,210,330}
      \draw[vec, blue, -{Stealth[length=6pt,width=5pt,bend]}]
        (\s-17:1.25) arc[start angle=\s-17, delta angle=-86, radius=1.25];
    \node[blue] at (0,0) {$-1$};
    \node[anchor=north, font=\footnotesize] at (0,-2.1)
      {antizyklisch: $132,\ 321,\ 213$};
  \end{scope}
\end{tikzpicture}
\end{document}
